\section{模型检验}

本文从数据、需求、选址容量、定价财务和情景重求解五类约束检验模型输出的一致性。财务指标沿用问题三式 \eqref{eq:q3-topic-profit-rate}--\eqref{eq:q3-annual-net} 的正式定义，本节不重复展开计算过程。检验结果汇总见表 \ref{tab:model-validation-summary}。

\begin{table}[H]
  \centering
  \caption{模型约束与一致性检验汇总}
  \label{tab:model-validation-summary}
  \small
  \renewcommand{\arraystretch}{1.18}
  \begin{tabular}{p{0.16\textwidth}p{0.51\textwidth}p{0.25\textwidth}}
    \toprule
    检验对象 & 检验内容 & 检验结论 \\
    \midrule
    数据完整性 &
    十个小区、三类老人、六项服务及距离矩阵均完整；距离矩阵对称且对角元为零；紧急救助基准价格为零。 &
    输入数据与公益服务口径一致。 \\
    需求约束 &
    预测人口与需求均非负；收费服务在消费预算不足时等比例削减；紧急救助需求保持理论值。 &
    收费需求满足预算约束，公益服务未被削减。 \\
    选址约束 &
    最优方案为 $C$ 小型、$D$ 大型和 $G$ 大型站，总建设成本 $108$ 万元；被分配小区至服务站距离均不超过 $1000$ 米。 &
    满足建设预算与服务半径约束。 \\
    容量约束 &
    三站容量可得系数分别为 $0.9978$、$0.9767$ 和 $1.0000$，利用率均不超过 $1$。 &
    容量折减后服务量可行。 \\
    定价与财务 &
    候选价格均反馈重算满意度、分配和财务指标；入选方案满足 $\rho_j^{topic}\leq 8\%$，且未要求 $\rho_j^{topic}\geq0$；$C$ 站利润率和年度净收益为负。 &
    利润率约束通过，但 $C$ 站存在保本风险。 \\
    情景一致性 &
    $S0$--$S4$ 均重新执行需求、选址和定价补贴求解；$S2$--$S4$ 的年度净收益为负。 &
    灵敏度结果可比，扰动情景存在财务风险。 \\
    \bottomrule
  \end{tabular}
\end{table}

综上，模型满足需求预算、服务半径、建设预算、容量和利润率上限约束，同时保留了对负利润率、负年度净收益及预算扩张下运营风险的识别。完整阶段检查记录见 \texttt{outputs/logs/}。
